Questo capitolo descrive il disegno sperimentale con cui la tesi verifica le proprie ipotesi: i dati e l'addestramento di STERA 30M, i modelli di riferimento, il protocollo di valutazione predittiva e narrativa, l'analisi statistica e le misure di memoria e tempo.

\section{Ipotesi di lavoro}
\label{sec:ipotesi}

Le domande di ricerca dell'Introduzione si traducono in quattro ipotesi verificabili, ciascuna associata a una misura e a un confronto.

\begin{table}[htbp]
\centering
\small
\begin{tabularx}{\linewidth}{lXX}
\toprule
& Ipotesi & Misura e confronto \\
\midrule
H1 & STERA 30M apprende regolarità linguistiche che generalizzano a racconti esclusi dall'addestramento. & NLL su DEV durante il training; NLL sui racconti TEST, confrontata con i due riferimenti. \\
H2 & La parametrizzazione STERA riduce i byte necessari per rappresentare i pesi. & Payload dei parametri compatti rispetto alla propria espansione densa e ai checkpoint di riferimento. \\
H3 & Applicando le proiezioni direttamente alle coordinate, la generazione è più rapida che con le matrici dense equivalenti, a parità di pesi. & Tempo per token in decode dei due motori sullo stesso modello, con verifica dell'identità dei token generati. \\
H4 & STERA 30M genera racconti di qualità confrontabile con quella dei riferimenti TinyStories. & Quota di racconti corretti in una valutazione cieca sugli stessi incipit, con confronti appaiati. \\
\bottomrule
\end{tabularx}
\caption{Ipotesi della tesi e misure con cui sono verificate.}
\label{tab:ipotesi}
\end{table}

\section{Corpus e separazione di TRAIN, DEV e TEST}\label{sec:stera30m-dati}

Il corpus di STERA 30M deriva dal file \texttt{TinyStoriesV2-GPT4-train} \cite{eldan2023tinystories}. I tre insiemi della \autoref{tab:stera30m-split} sono partizioni interne di quel file. Le storie sono tokenizzate con un vocabolario BPE di \num{50257} voci. L'identificativo 50.256 delimita l'inizio e la fine di ciascun racconto.

\begin{table}[htbp]
\centering
\small
\begin{tabularx}{\linewidth}{lrrX}
\toprule
Insieme & Racconti & Token target & Uso \\
\midrule
TRAIN & \num{2660865} & \num{530189191} & Aggiornamento dei parametri \\
DEV & \num{29279} & \num{5838521} & Controllo durante l'addestramento e valutazione su un campione distinto \\
TEST & \num{27351} & \num{5441693} & Prove predittive e narrative dopo il congelamento di STERA 30M \\
\bottomrule
\end{tabularx}
\caption{Partizioni del corpus di STERA 30M. I token target contano le posizioni predicibili di ogni racconto, compreso il token di fine.}
\label{tab:stera30m-split}
\end{table}

La separazione avviene per racconto. Il programma uniforma gli spazi bianchi e converte in minuscolo i byte ASCII prima di calcolare SHA-256. Interpreta i primi otto byte dell'impronta come intero in ordine \emph{little-endian} e ne prende il resto modulo \num{10000}. I valori da 0 a 99 vanno in TEST, quelli da 100 a 199 in DEV e gli altri in TRAIN. I \num{2000} gruppi di contenuto già riservati a una valutazione separata sono assegnati a DEV. Storie identiche dopo la normalizzazione ricevono la stessa assegnazione. L'audit ha verificato tutte le righe dell'indice. Testi simili ma con impronte diverse possono trovarsi in partizioni diverse.

\section{Configurazione del pretraining}
\label{sec:stera30m-config-training}

Il modello parte da zero ed esegue \num{122000} aggiornamenti, per un totale di \num{1529514422} token target: circa $2{,}9$ passaggi sull'insieme TRAIN. Ogni aggiornamento usa un batch di 64 sequenze, suddiviso fra due GPU NVIDIA P100, con accumulo dei gradienti su gruppi di otto sequenze per GPU. Le finestre arrivano a 512 token target e restano entro lo stesso racconto; il padding è escluso dalla perdita. L'ordine di visita di TRAIN è deterministico: il primo indice è \num{194723} e il passo è \num{104729} modulo il numero di racconti di TRAIN.

AdamW usa $\beta_1=0{,}9$, $\beta_2=0{,}95$ e limita la norma del gradiente a 1. Il weight decay è $0{,}1$ per coordinate RowMask23, fattori a basso rango e vettori di posizione. Il learning rate segue un profilo a tre fasi (\autoref{fig:stera30m-pretraining}): sale linearmente fino a $6\cdot10^{-4}$ nei primi 300 passi, resta costante fino al passo \num{64000}, scende linearmente fino a $6\cdot10^{-5}$ al passo \num{80000} e rimane a questo valore fino al termine. La NLL su DEV è misurata prima del training e ogni \num{2000} aggiornamenti, sempre sui primi 512 racconti DEV nell'ordine dell'indice (\num{102820} token target, al massimo 512 per racconto). Il modello valutato nei capitoli successivi è il checkpoint finale.

\begin{table}[htbp]
\centering
\small
\begin{tabularx}{\linewidth}{Xl}
\toprule
Operazione & Seme \\
\midrule
Seme globale PyTorch & \texttt{20261117} \\
Inizializzazione delle coordinate e dei vettori di posizione & \texttt{20261121} \\
Inizializzazione del fattore a basso rango $V$ & \texttt{7} \\
Selezione dei 300 incipit TEST & \texttt{202611240901} \\
Campionamento della continuazione $i$, con $i$ da 0 a 299 & $202611240902+1009\,i$ \\
\bottomrule
\end{tabularx}
\caption{Semi del pretraining e della valutazione. Il fattore $U$ parte da zero. Lo split è determinato dalla regola di hash della \autoref{sec:stera30m-dati}.}
\label{tab:stera30m-seed}
\end{table}


\section{Modelli di riferimento}
\label{sec:metodo-modelli-tre}

Il confronto linguistico impiega STERA 30M e due checkpoint pubblici della famiglia TinyStories \cite{eldan2023tinystories}, addestrati sullo stesso dominio e con lo stesso tokenizzatore (\autoref{tab:modelli-principali}). Il numero nel nome pubblico dei checkpoint TinyStories non coincide con il conteggio dei loro parametri: per tutti i modelli riportiamo il numero effettivo di valori memorizzati, contando una sola volta le matrici condivise.

\begin{table}[htbp]
\centering
\small
\begin{tabularx}{\linewidth}{Xrrr}
\toprule
Modello & Parametri & Pesi FP32 (MB) & Strati \\
\midrule
TinyStories-8M & \num{19702528} & 78,810 & 8 \\
\textbf{STERA 30M} & \textbf{\num{30346182}} & \textbf{121,385} & 6 \\
TinyStories-33M & \num{68514048} & 274,056 & 4 \\
\bottomrule
\end{tabularx}
\caption{Modelli del confronto linguistico. MB indica milioni di byte e comprende i soli parametri in FP32.}
\label{tab:modelli-principali}
\end{table}

STERA 30M si colloca fra i due riferimenti per numero di parametri memorizzati: ne ha circa una volta e mezza quelli di TinyStories-8M e meno della metà di quelli di TinyStories-33M. Durante il calcolo, invece, le sue matrici espanse contano \num{82721280} valori: il costo aritmetico per token di STERA 30M è dunque paragonabile a quello di un modello denso di dimensione simile a TinyStories-33M, e il vantaggio atteso riguarda la memoria, non il numero di operazioni.

\section{Protocollo di valutazione}
\label{sec:metodo-confronto-generativo}

La \autoref{fig:protocollo} riassume il flusso completo, dai dati al confronto statistico.

\begin{figure}[htbp]
\centering
\resizebox{\linewidth}{!}{\begin{tikzpicture}[font=\footnotesize, >={Latex[length=2mm]},
  st/.style={draw=steraGrigio, rounded corners=2pt, fill=white, minimum width=30mm, minimum height=12mm, align=center, text width=28mm},
  hi/.style={st, draw=steraBlu, fill=steraBlu!10}]
\node[st] (a) {Corpus TinyStories,\\split per hash};
\node[st, right=6mm of a] (b) {Training su TRAIN, controllo su DEV};
\node[st, right=6mm of b] (c) {Pesi congelati};
\node[st, right=6mm of c] (d) {300 prompt da TEST, seme fisso};
\node[hi, below=10mm of d] (e) {Generazione dei 3 modelli, stesso campionamento};
\node[hi, left=6mm of e] (f) {900 testi anonimi, ordine casuale};
\node[hi, left=6mm of f] (g) {Giudizio cieco con rubrica};
\node[hi, left=6mm of g] (h) {Abbinamento ai modelli e statistica};
\draw[->] (a) -- (b); \draw[->] (b) -- (c); \draw[->] (c) -- (d); \draw[->] (d) -- (e);
\draw[->] (e) -- (f); \draw[->] (f) -- (g); \draw[->] (g) -- (h);
\node[st, below=10mm of e, fill=steraArancio!10, draw=steraArancio] (nll) {NLL in teacher forcing sui 300 racconti};
\draw[->] (d.east) -- ++(4mm,0) |- (nll.east);
\end{tikzpicture}
}
\caption{Protocollo di validazione. In alto la preparazione del modello e degli incipit, in basso la valutazione: la NLL sui racconti di riferimento e il giudizio cieco sulle continuazioni generate.}
\label{fig:protocollo}
\end{figure}

\subsection{Incipit di prova}

Dalla partizione TEST sono estratti, con un seme fissato, 300 racconti. Di ciascuno si conservano come incipit le prime due frasi complete, purché contengano fra 12 e 64 token e lascino almeno 20 token di continuazione. La selezione usa solo il testo di riferimento e non dipende da alcun modello. Gli stessi 300 incipit sono usati per tutte le misure di questo capitolo.

\subsection{Verosimiglianza dei racconti di riferimento}

Per ogni modello si calcola la NLL dei 300 racconti TEST completi in \emph{teacher forcing}: a ogni posizione il modello riceve il prefisso reale e si misura la probabilità assegnata al token successivo corretto. I racconti contano complessivamente \num{57841} token target, con un massimo di 512 per racconto. La NLL è la media ponderata per token e la perplessità il suo esponenziale. I tre modelli condividono lo stesso tokenizzatore, quindi le NLL sono direttamente confrontabili.

\subsection{Generazione delle continuazioni}

Ogni modello continua ciascun incipit una sola volta, con la stessa procedura: campionamento top-$k$ con $k=40$ e temperatura $0{,}75$, al più 180 nuovi token, arresto al token di fine racconto, pesi in FP32 e seme del generatore dipendente dall'indice dell'incipit (\autoref{tab:stera30m-seed}). Poiché il campionamento usa un generatore su CPU con lo stesso seme per tutti i modelli, le differenze fra le continuazioni dipendono solo dalle distribuzioni prodotte dai modelli.

\subsection{Giudizio cieco}
\label{sec:metodo-giudizio}

Le 900 continuazioni dei tre modelli sono rese anonime, mescolate in ordine casuale e divise in blocchi da 30. Ogni testo riceve un identificativo che non contiene informazioni sulla provenienza; la corrispondenza fra identificativi e modelli è conservata separatamente e usata solo dopo la conclusione del giudizio. I testi sono valutati da un modello linguistico, Claude Opus 5.5 \cite{anthropic2026claude}, uno alla volta, con la rubrica riportata in Appendice e ispirata a quella di Eldan e Li \cite{eldan2023tinystories}. Per ogni testo il giudice assegna tre voti interi da 1 a 10:
\begin{itemize}
\item \textbf{grammatica}: correttezza e scorrevolezza della lingua;
\item \textbf{coerenza}: stabilità di identità, attributi, oggetti e fatti, anche rispetto all'incipit;
\item \textbf{trama}: presenza di eventi collegati da una causalità comprensibile;
\end{itemize}
e una categoria: S (coerente), P (nucleo riconoscibile ma almeno una criticità sostanziale, con almeno un voto non superiore a 5) o N (perdita grave di coerenza, con almeno un voto non superiore a 3). Ogni giudizio è accompagnato da una motivazione specifica. Un racconto è \emph{corretto} se è classificato S, cioè se tutti e tre i voti sono almeno 6:
\begin{equation}
I_i=\mathbf 1\{g_i\geq6,\ c_i\geq6,\ p_i\geq6\}.
\label{eq:regola-correttezza}
\end{equation}
Fantasia, animali parlanti, finali negativi e assenza di una morale non sono penalizzati; lo sono le contraddizioni, i cambi immotivati di identità, le ripetizioni prive di significato e l'incoerenza causale. Una continuazione interrotta dal limite di lunghezza non è considerata un errore se il nucleo narrativo è comprensibile.

\section{Analisi statistica}
\label{sec:metodo-statistica}

\begin{itemize}
\item \textbf{Proporzioni.} La quota di racconti corretti di ciascun modello è accompagnata dall'intervallo di Wilson al 95\%.
\item \textbf{Differenze appaiate.} Per ogni coppia di modelli la differenza di correttezza è stimata sugli stessi 300 incipit; l'intervallo al 95\% è ottenuto con un bootstrap che ricampiona gli incipit tenendo unite le valutazioni dei modelli, con \num{25000} repliche.
\item \textbf{Significatività.} La differenza è verificata con il test esatto di McNemar sulle coppie discordanti; i tre confronti a coppie sono corretti con il metodo di Holm.
\item \textbf{NLL.} La differenza di NLL fra due modelli è accompagnata da un intervallo bootstrap al 95\% che ricampiona i racconti, con \num{10000} repliche.
\end{itemize}

\section{Misure di memoria e tempo}
\label{sec:metodo-sistemi}

Le misure distinguono il payload dei parametri, la memoria fisica occupata dal processo, il tempo al primo token e il tempo per i token successivi. Il confronto principale riguarda due motori che eseguono gli stessi pesi di STERA 30M: uno applica le proiezioni direttamente alle coordinate RowMask23, l'altro carica le matrici dense equivalenti. I tempi di esecuzione dipendono dalla forma del modello e non dai valori dei pesi. Lo studio dell'esecuzione con pesi che superano la memoria disponibile, descritto nel \autoref{cap:capitoli/07-sistemi}, usa una configurazione STERA di scala maggiore e un modello pubblico eseguito con caricamento progressivo dei pesi dal disco.

\paragraph{Protocollo delle misure di tempo.} Tutte le misure di tempo sono eseguite sullo stesso Mac con chip M1 e 16~GiB di memoria unificata, con i pesi del checkpoint finale già caricati, batch 1, greedy e cache KV di capacità 512. Ogni configurazione è misurata in cinque campagne complete; ciascuna campagna parte solo quando la media di carico del sistema a un minuto è scesa sotto 1,5, dopo una pausa di 60 secondi. Ogni processo esegue una prima richiesta e una seconda a caldo, che è quella riportata; i valori sono mediane fra campagne e input, e in ogni campagna si verifica che i token generati coincidano fra le varianti e fra le campagne.

\paragraph{Confronto alla pari con i riferimenti.} I checkpoint TinyStories sono eseguiti in due modi. Il primo è il motore standard con cui sono distribuiti, Transformers su PyTorch con backend MPS. Il secondo è un motore denso MLX con le stesse ottimizzazioni del motore denso di STERA 30M, che carica i pesi pubblici e produce gli stessi token dell'implementazione di riferimento. Solo il secondo confronto attribuisce le differenze di tempo ai modelli anziché ai framework.

\paragraph{Variazioni dell'ingresso e del prefill.} Oltre ai 256 token di ingresso delle misure principali, si misurano prefissi di 32, 64, 128 e 384 token, sempre con 100 token generati. Per il motore compatto si misurano due modi di elaborare il prefisso: il prodotto diretto a blocchi di token e l'espansione al bisogno delle matrici strato per strato.

\section{Analisi delle componenti delle proiezioni}
\label{sec:metodo-componenti}

Per misurare quanto il modello addestrato dipenda da scale, guadagni e termine a basso rango, e se queste componenti si possano rappresentare con meno bit, si applicano a posteriori, sul checkpoint finale, le seguenti modifiche:
\begin{itemize}
\item \textbf{eliminazione} di scale o guadagni, riportati a 1;
\item \textbf{riduzione del rango} del termine $UV^\top$ alla migliore approssimazione di rango $k\in\{64,48,32,16,8,0\}$, ottenuta con la decomposizione ai valori singolari;
\item \textbf{quantizzazione a 8 bit} di scale, guadagni o fattori $U$ e $V$, simmetrica con una scala per matrice per scale e guadagni e per gruppi di valori per $U$ e $V$; per i fattori si misura anche l'errore relativo della matrice densa equivalente.
\end{itemize}
Per ogni variante si calcola la NLL sui 300 racconti TEST. Per le varianti del basso rango si generano anche le continuazioni dei 300 incipit, con la stessa procedura e gli stessi semi, e si ripete il giudizio cieco: i testi di tutte le varianti sono mescolati e anonimizzati e valutati dallo stesso giudice con la stessa rubrica, in sessioni separate da quella dei tre modelli. Queste prove misurano la dipendenza del modello già addestrato dalle singole componenti, non l'effetto di addestrare il modello senza di esse.
